\chapter{仿射空间}\label{ch:b2:affine}

向量空间有一个特殊的点——原点——
几何并不需要它。一个\emph{\omterm{def:b2:affine:def}{仿射空间}}是
忘记了其原点的向量空间：点与向量成为
两类不同的对象，由平移相互联系。本章
建立这一“字典”（点、\omterm{def:b2:affine:barycenter}{重心}、仿射子空间与\omterm{def:b2:affine:subspace}{仿射映射}），
凸性的仿射观点，以及
后续几何章节中用到的分类工具。

\section{点与向量}

\begin{definition}\label{def:b2:affine:def}
以实向量空间 $E$ 为方向的\emph{仿射空间}\index{仿射空间}是
一个非空集合 $\mathcal{E}$，并配有映射
$(A, B) \mapsto \vect{AB} \in E$，满足
\[
\vect{AB} + \vect{BC} = \vect{AC}
\quad \text{（沙勒关系）},
\qquad
\text{对每个 } A,\ B \mapsto \vect{AB} \text{ 是双射 }
\mathcal{E} \to E .
\]
记 $B = A + u$ 为满足 $\vect{AB} = u$ 的唯一点。
\emph{维数}：$\mathcal{E}$ 的维数为 $\dim E$。每个向量
空间都是其自身上的仿射空间（$\vect{AB} = B - A$）；
每一种原点选取 $O \in \mathcal{E}$ 都将 $\mathcal{E}$ 与
$E$ 通过 $M \mapsto \vect{OM}$ 等同起来。
\end{definition}

\begin{example}[没有自然
原点的仿射空间]\label{ex:b2:affine:planeexample}
解平面 $\mathcal E = \{(x, y, z) \in \R^3 : x + y
+ z = 1\}$ 不是向量子空间（$0 \notin \mathcal E$），
但它是以 $E = \{x + y + z =
0\}$ 为方向的\omterm{def:b2:affine:def}{仿射空间}：对 $A, B \in \mathcal E$，差 $\vect{AB} = B
- A$ 落在 $E$ 中（和式相消），沙勒关系
继承自 $\R^3$，且 $B \mapsto \vect{AB}$ 是到 $E$ 上的双射。
$\mathcal E$ 中没有哪一点被区分出来——任取
“原点”$O \in \mathcal E$ 同样可用，且所有的
等同 $M \mapsto \vect{OM}$ 彼此只差一个平移。
这是典型情形：非齐次线性问题的解集
（线性方程组、线性\omterm{def:b2:diffcalc:differential}{微分}
方程，见\cref{ch:b2:diffeq}）都是仿射的，从来不是线性的，
而“特解加核”这一口号恰好就是
下一定义中的陈述 $\mathcal F = A + F$。
\end{example}

\begin{definition}[重心]\label{def:b2:affine:barycenter}
设 $(A_i, \lambda_i)_{i \leq k}$ 为带权点，且
$\sum\lambda_i \neq 0$。称\emph{重心}\index{重心} $G
= \operatorname{bar}\bigl((A_i, \lambda_i)\bigr)$ 为满足
下式的唯一点：
\[
\sum_i \lambda_i\, \vect{GA_i} = 0
\qquad\text{等价地}\qquad
\vect{OG} = \frac{1}{\sum\lambda_i}\sum_i \lambda_i\,\vect{OA_i}
\quad (\text{对任意 } O).
\]
重心具有\emph{结合性}（点的子族可
代之以其部分重心，权重取为其和）；且
在所有权重同乘以一个非零数时保持不变。
\end{definition}

\begin{proof}[存在性与公式的证明]
取定 $O$，记 $s = \sum_i\lambda_i \neq 0$。由沙勒关系，
\[
\sum_i\lambda_i\,\vect{GA_i} = 0
\iff \sum_i\lambda_i\bigl(\vect{GO} + \vect{OA_i}\bigr) = 0
\iff s\,\vect{OG} = \sum_i\lambda_i\,\vect{OA_i},
\]
由此 $G = O + \frac1s\sum_i\lambda_i\vect{OA_i}$
唯一确定。\emph{与 $O$ 无关：}对另一原点
$O'$，
\[
\frac1s\sum_i\lambda_i\,\vect{O'A_i}
= \frac1s\sum_i\lambda_i\bigl(\vect{O'O} + \vect{OA_i}\bigr)
= \vect{O'O} + \frac1s\sum_i\lambda_i\,\vect{OA_i}
= \vect{O'G} :
\]
即同一点 $G$。\emph{结合性：}将指标集
分为 $I \sqcup J$，其中 $s_I = \sum_{i\in I}\lambda_i \neq 0$，
并记 $G_I$ 为 $(A_i, \lambda_i)_{i\in
I}$ 的\omterm{def:b2:affine:barycenter}{重心}，从而 $\sum_{i\in I}\lambda_i\vect{OA_i} =
s_I\,\vect{OG_I}$。于是
\[
s\,\vect{OG}
= \sum_{i\in I}\lambda_i\vect{OA_i} +
\sum_{j\in J}\lambda_j\vect{OA_j}
= s_I\,\vect{OG_I} + \sum_{j\in J}\lambda_j\vect{OA_j} :
\]
$G$ 即为 $(G_I, s_I)$ 与 $(A_j,
\lambda_j)_{j\in J}$ 的\omterm{def:b2:affine:barycenter}{重心}，正如所断言。\emph{伸缩：}将
每个 $\lambda_i$ 换为 $t\lambda_i$（$t \neq 0$），使 $s$
与加权和都乘以 $t$，从而 $\vect{OG}$ 保持不变。
\end{proof}

\begin{remark}[常见误区]
围绕本定义有两个陷阱。首先，若权重之和
为\emph{零}，则不存在\omterm{def:b2:affine:barycenter}{重心}：映射 $O \mapsto
\sum\lambda_i\vect{OA_i}$ 此时与 $O$ 无关，并
确定一个\emph{向量}，而非点——例如
$(A, -1;\ B, 1)$ 表示 $\vect{AB}$。时刻留意
计算得到的究竟是两者中的哪一个，是
\omterm{def:b2:affine:barycenter}{重心}计算基本功的一半。其次，权重只在
相差一个公共非零因子时才有确定意义；形如“
$G$ 的坐标为 $\lambda_1, \dots, \lambda_k$”的说法预设了
一种归一化（通常取 $\sum\lambda_i = 1$），而忘记
归一化，正是
图上比值出错的常见根源。
\end{remark}

\begin{definition}[仿射子空间；仿射映射]\label{def:b2:affine:subspace}
称\emph{仿射子空间}为集合 $\mathcal{F} = A + F = \{A + u :
u \in F\}$，$F$ 为向量子空间（其\emph{方向}）；
等价地，即对\omterm{def:b2:affine:barycenter}{重心}运算封闭的非空集合。
$\R^n$ 中的仿射子空间恰为线性方程组
$MX = B$ 的解集（本科一年级：特解加上核）。设映射 $f \colon
\mathcal{E} \to \mathcal{E}'$ 称为\emph{仿射映射}\index{仿射映射}，
若它保持\omterm{def:b2:affine:barycenter}{重心}——等价地，当
\[
f(A + u) = f(A) + \varphi(u)
\]
对某个（唯一的）线性映射 $\varphi = \vec f$，称其为\emph{线性
部分}。$\R^n$ 上的仿射映射即 $X \mapsto MX + C$。复合映射
仍为仿射映射，线性部分为各线性部分的复合；$f$ 为双射当且仅当 $\vec f$
为双射。
\end{definition}

\begin{proof}[映射等价性的证明]
若 $f(A + u) = f(A) + \varphi(u)$：设 $G$ 为 $(A_i,
\lambda_i)$ 的\omterm{def:b2:affine:barycenter}{重心}，以 $A$ 为起点展开各点，$f(G) = f(A) +
\varphi(\vect{AG})$ 且 $\varphi(\vect{AG}) =
\frac{\sum\lambda_i\varphi(\vect{AA_i})}{\sum\lambda_i}$，故 $f(G)$
为各像点的\omterm{def:b2:affine:barycenter}{重心}。反之，固定 $A$，定义
$\varphi(u) = \vect{f(A)\,f(A + u)}$。\emph{齐次性：}$A + tu =
\operatorname{bar}\bigl(A, 1-t;\ A + u, t\bigr)$，对每个
实数 $t$，故由所设的保持\omterm{def:b2:affine:barycenter}{重心}性质（权重取任意实数）
直接得 $\varphi(tu) = t\,\varphi(u)$。
\emph{可加性：}$A + u + v = \operatorname{bar}\bigl(A + 2u,
\tfrac12;\ A + 2v, \tfrac12\bigr)$，于是 $\varphi(u + v) =
\tfrac12\varphi(2u) + \tfrac12\varphi(2v) = \varphi(u) +
\varphi(v)$，利用齐次性，故 $\varphi$ 为线性映射。
\end{proof}

\begin{remark}
证明中所用的\omterm{def:b2:affine:barycenter}{重心}权重为\emph{任意实数}，
即齐次性一步中 $t$ 可取 $\intcc01$ 之外的值。
若仅假设映射保持权重非负的\omterm{def:b2:affine:barycenter}{重心}——
等价地，即保持中点与线段，
则向量映射的线性性不再自动成立：
仅能得到 $\Q$-线性，要得出结论还需
一个\omterm{def:b2:metric:continuity}{连续}性假设，正如
\cref{exo:b2:affine:5}中那样。区分“保持所有
\omterm{def:b2:affine:barycenter}{重心}”与“保持凸组合”是仿射语言中的
一个虽小却真实的细微之处。
\end{remark}

\begin{example}[经典的重心几何]\label{ex:b2:affine:median}
三角形 $ABC$ 的形心即\omterm{def:b2:affine:barycenter}{重心} $G =
\operatorname{bar}(A,1; B,1; C,1)$。由结合律及
中点 $A' = \operatorname{bar}(B, 1; C, 1)$，可得
\[
G = \operatorname{bar}(A, 1;\ A', 2) :
\]
$G$ 位于中线 $AA'$ 上，且在其三分之二处——另外两条中线同理，
三条中线共点，这在\omterm{def:b2:affine:barycenter}{重心}演算中
只需一行即可看出。
\end{example}

\begin{example}[四边形的
双中线]\label{ex:b2:affine:bimedians}
设 $ABCD$ 为任一四边形（平面的或非平面的！），并考虑
其\emph{双中线}，即连接
对边中点的线段 $M_{AB}M_{CD}$ 与 $M_{BC}M_{DA}$。引入
$G$ 为 $(A,1; B,1; C,1; D,1)$ 的\omterm{def:b2:affine:barycenter}{重心}，并将权重按
两种方式分组：
\[
G = \operatorname{bar}\bigl(M_{AB}, 2;\ M_{CD}, 2\bigr)
= \operatorname{bar}\bigl(M_{BC}, 2;\ M_{DA}, 2\bigr) :
\]
$G$ 是\emph{两}条双中线的中点——故两条双中线
总互相平分，四个中点所成的四边形
为平行四边形（其对角线即双中线）。
无需分类讨论，也无需坐标，论证
原样适用于
$\R^3$ 中的斜四边形，图形论证在此已颇为棘手：
结合律与维数无关。
\end{example}

\begin{example}[仿射映射的分类：从头
到尾]\label{ex:b2:affine:dilation}
设 $f(x, y) = (2x - 1,\ 3y - 4)$，定义在 $\R^2$ 上。其线性部分
为 $\varphi = \operatorname{diag}(2, 3)$，其谱
$\{2, 3\}$ 不含 $1$：由下文证明的不动点判别法
（见\cref{prop:b2:affine:fixedpoint}），$f$ 恰有一个
不动点，其可通过解下列方程求得：
\[
x = 2x - 1, \qquad y = 3y - 4
\qquad\Longrightarrow\qquad \Omega = (1, 2).
\]
重新以 $\Omega$ 为原点（取 $x = 1 + u$，$y = 2 + v$）：
\[
f(1 + u,\ 2 + v) = (1 + 2u,\ 2 + 3v) :
\]
在以 $\Omega$ 为原点的坐标系中，$f$ \emph{正是}其线性部分，
一个各向异性伸缩，水平方向拉伸 $2$ 倍，
竖直方向拉伸 $3$ 倍，中心为 $(1, 2)$。一般的道理是：
\omterm{def:b2:affine:subspace}{仿射映射}是“线性映射加上定位数据”，而
定位数据可归结为一个选得恰当的原点，只要
$1$ 不是\omterm{def:b2:reduction:eigen}{特征值}。反之，平移原点不当，
便会\emph{产生}常数项：仿射几何
即选定把 $0$ 放在何处的艺术。
\end{example}

\begin{remark}[方法：借重心证明
共点与共线]
\cref{ex:b2:affine:median}是一般方法的一个实例。
要证三角形的三条塞瓦线共点，
给出\emph{单一}的加权系统 $(A,
\alpha; B, \beta; C, \gamma)$，并从三种方式运用结合律：
分组 $(B, C)$ 可知\omterm{def:b2:affine:barycenter}{重心}位于由
$A$ 引出的塞瓦线上；分组 $(C, A)$ 则位于由 $B$ 引出的
塞瓦线上，分组 $(A, B)$ 则位于第三条。对中线，
系统 $(A, 1; B, 1; C, 1)$ 即完成全部工作；对于
截分各边为指定比的塞瓦线，其权重由所给的
比值直接读出。要证三点\emph{共线}，将其中一点写成
另外两点的\omterm{def:b2:affine:barycenter}{重心}（见\cref{exo:b2:affine:2}），或
用\cref{exo:b2:affine:11}中的\omterm{def:b2:linalg:det}{行列式}判别法。两种
方法都以权重的推算取代几何上的巧思——
这正是\omterm{def:b2:affine:barycenter}{重心}演算的用武之地。
\end{remark}

\section{仿射凸性}

\begin{definition}\label{def:b2:affine:convex}
\omterm{def:b2:affine:def}{仿射空间}的子集 $C$ 称为\emph{凸集}\index{凸集}，当它包含
其各点以\emph{非负}权重所得的每个\omterm{def:b2:affine:barycenter}{重心}——
等价地，其中任意两点间的线段 $\intcc{A}{B} = \{\operatorname{bar}(A,
1-t; B, t) : t \in \intcc{0}{1}\}$ 皆含于该集。
\emph{凸包}\index{凸包} $\operatorname{conv}(S)$ 是由 $S$ 中各点的
全部非负权重\omterm{def:b2:affine:barycenter}{重心}所组成的集合——即
包含 $S$ 的最小凸集。
\end{definition}

\begin{example}[上图为
凸集]\label{ex:b2:affine:epigraph}
区域 $C = \{(x, y) : y \geq x^2\}$ 位于抛物线
上方，且为\omterm{def:b2:affine:convex}{凸集}：对 $(x_1, y_1), (x_2, y_2) \in C$ 及 $t \in
\intcc01$，由平方函数的凸性不等式
可得
\[
\bigl((1-t)x_1 + tx_2\bigr)^2 \leq (1-t)x_1^2 + tx_2^2
\leq (1-t)y_1 + ty_2 ,
\]
故该\omterm{def:b2:affine:barycenter}{重心}仍位于抛物线上方。此计算
具有一般性：$\{y \geq f(x)\}$ 为\omterm{def:b2:affine:convex}{凸集}，当且仅当 $f$ 为
凸函数——凸\emph{集}与凸
\emph{函数}（见\cref{ch:b2:realfun}）是同一概念的两个侧面，
上图即二者之间的对应。这正是
凸函数存在支撑线的几何原因，此事实
将在
\cref{ch:b2:randomvar}中用于证明琴生不等式。
\end{example}

\begin{example}[凸包的
冗余生成元]\label{ex:b2:affine:squarehull}
设 $S = \{(0,0), (2,0), (2,2), (0,2), (1,1)\}$。第五个
点即为\omterm{def:b2:affine:barycenter}{重心}：
\[
(1,1) = \operatorname{bar}\bigl((0,0), \tfrac12;\ (2,2),
\tfrac12\bigr),
\]
故它已位于其余四点的\omterm{def:b2:affine:convex}{凸包}之中，
$\operatorname{conv}(S)$ 即是以四个角点
为顶点的正方形。一般地，$S$ 中某点若为
\emph{其余}各点的非负权重\omterm{def:b2:affine:barycenter}{重心}，
则可删去该点，而 $S$ 的\omterm{def:b2:affine:convex}{凸包}保持不变；不能删去的点
（此处即四个角点）为\omterm{def:b2:affine:convex}{凸包}的
\emph{极点}。确定它们纯属
\omterm{def:b2:affine:barycenter}{重心}计算：以 $(2,0)$ 为例，它不能写成其余各点
以非负权重所作的 $\operatorname{bar}$，因第一坐标
会迫使全部权重落在
$x = 2$ 的点上，于是第二坐标随即不成立。凸性问题
一再归结为
求解若干小型加权方程组。
\end{example}

\begin{theorem}[卡拉泰奥多里]\label{thm:b2:affine:caratheodory}
在维数为 $n$ 的\omterm{def:b2:affine:def}{仿射空间}中，
$\operatorname{conv}(S)$ 的每一点都是至多 $n + 1$ 个
$S$ 中点的\omterm{def:b2:affine:barycenter}{重心}。\index{卡拉泰奥多里定理}
\end{theorem}

\begin{proof}
设 $G = \operatorname{bar}(A_0, \lambda_0; \dots; A_k, \lambda_k)$
其中 $\lambda_i > 0$，$\sum\lambda_i = 1$，且 $k + 1 > n + 1$
个点。$k$ 个向量 $\vect{A_0A_i}$（$i \geq 1$）线性相关
（$k > n$），即有非平凡的 $\sum_{i\geq1}\mu_i \vect{A_0A_i} = 0$；
令 $\mu_0 = -\sum_{i\geq1}\mu_i$，于是得权重 $(\mu_i)$，
满足 $\sum\mu_i = 0$，$\sum \mu_i\,\vect{OA_i} = 0$（$O$ 任意），
且不全为零。于是对任意实数 $t$，权重 $\lambda_i - t\mu_i$
其和仍为 $1$，且由 $\sum_i\mu_i\vect{OA_i} = 0$，
\[
\sum_i(\lambda_i - t\mu_i)\,\vect{OA_i}
= \sum_i\lambda_i\,\vect{OA_i} :
\]
它们给出的是\emph{同一}点 $G$。现令 $t$ 自
$0$ 起增大：某个 $\mu_i$ 为正（其和为零，且
不全为零），故
\[
t^* = \min\Bigl\{\frac{\lambda_i}{\mu_i} : \mu_i > 0\Bigr\}
\]
该值良定且为正。当 $t = t^*$ 时，对满足
$\mu_i > 0$ 的指标，由极小性 $\lambda_i - t^*\mu_i \geq 0$，
且在取极小的指标处等号成立；对满足 $\mu_i
\leq 0$ 的指标，$\lambda_i - t^*\mu_i \geq \lambda_i > 0$。全部
权重仍皆非负，且至少一个已变为零，于是 $G$
可改写为点数更少的\omterm{def:b2:affine:barycenter}{重心}。反复此过程，
直至剩余点数不超过 $n + 1$ 个。
\end{proof}

\begin{example}
在平面（$n = 2$）中：有限集的\omterm{def:b2:affine:convex}{凸包}中每一点
都位于以该集中点为顶点的三角形之内——这正是
卡拉泰奥多里定理的几何内容，它在优化与
概率（混合分布）中皆有应用。
\end{example}

\begin{example}[卡拉泰奥多里
算法演示]\label{ex:b2:affine:caratheodoryrun}
记\cref{ex:b2:affine:squarehull}中正方形的中心，
其四个角点为 $A_1 =
(0,0)$, $A_2 = (2,0)$, $A_3 = (2,2)$, $A_4 = (0,2)$:
\[
(1,1) = \operatorname{bar}\bigl(A_1, \tfrac14;\ A_2,
\tfrac14;\ A_3, \tfrac14;\ A_4, \tfrac14\bigr),
\]
维数为 $2$ 的空间中有四个点，多出一个。证明中的
步骤要求权重 $(\mu_i)$ 满足 $\sum\mu_i = 0$ 且
$\sum\mu_i\vect{OA_i} = 0$：此处取 $\mu = (1, -1, 1, -1)$ 即可
（两条对角线共享中点）。令 $\lambda_i
\mapsto \lambda_i - t\mu_i$ 时，\omterm{def:b2:affine:barycenter}{重心}对每个 $t$ 保持不变；
可行的极值 $t = \frac14$ 使
权重成为 $(0, \tfrac12, 0, \tfrac12)$，同时消去 $A_1$ 与
$A_3$：
\[
(1,1) = \operatorname{bar}\bigl(A_2, \tfrac12;\ A_4,
\tfrac12\bigr),
\]
这是用两点表示的\omterm{def:b2:affine:barycenter}{重心}——比定理所保证的
三点还要好，原因在于中心恰好
位于两个\omterm{def:b2:structures:generated}{生成}元之间的线段上。此算法
完全是机械的：求出一个线性关系，沿参数滑动，直至某个权重
变为零，然后重复。
\end{example}

\section{仿射分类工具}

\begin{proposition}[仿射映射的不动点]\label{prop:b2:affine:fixedpoint}
设 $f$ 为有限维\omterm{def:b2:affine:def}{仿射
空间}上的仿射自同态，其线性部分为 $\varphi$。若 $1 \notin
\operatorname{Sp}(\varphi)$，则 $f$ 恰有一个不动点
$\Omega$，且在 $\Omega$ 处的向量化中，$f$ \emph{即为}
其线性部分。（平移，$\varphi = \mathrm{id}$ 且无
不动点，是最基本的障碍。）
\end{proposition}

\begin{proof}
取定 $O$，并记 $f(O + x) = f(O) + \varphi(x)$。点
$O + x$ 为不动点当且仅当 $O + x = f(O) + \varphi(x)$，即
\[
(\mathrm{id} - \varphi)(x) = \vect{O f(O)} .
\]
有限维时，$\mathrm{id} - \varphi$ 可逆
当且仅当 $0$ 不是 $\mathrm{id} - \varphi$ 的\omterm{def:b2:reduction:eigen}{特征值}，亦即
$1 \notin \operatorname{Sp}\varphi$——此时
上式恰有一解 $x^*$，给出
唯一不动点 $\Omega = O + x^*$。重新取原点：对任意
向量 $u$，
\[
f(\Omega + u) = f(\Omega) + \varphi(u) = \Omega + \varphi(u),
\]
于是在以 $\Omega$ 为原点的坐标系中，该映射为 $u \mapsto
\varphi(u)$，即纯线性的。当 $1 \in
\operatorname{Sp}\varphi$ 时，要么不存在不动点（上式
可能无解，如平移的情形），要么
整个仿射子空间都是不动点（在某一解上加任一
\omterm{def:b2:reduction:eigen}{特征值}为 $1$ 的\omterm{def:b2:reduction:eigen}{特征向量}即可），唯一性恰好
由谱条件刻画。
\end{proof}

\begin{example}[平面等距变换的完整分类]\label{ex:b2:affine:isometries}
欧几里得平面上的仿射等距变换，其线性部分属于
$O(2)$：旋转 $R_\theta$ 或反射（参见本科一年级分册）。若
$\theta \neq 0$，则 $1 \notin \operatorname{Sp} R_\theta$，因此该映射
是\emph{绕唯一中心的旋转}
（见\cref{prop:b2:affine:fixedpoint}）。若线性部分
为反射，则或为关于某轴的反射（存在不动点），或为
\emph{滑动反射}（反射与沿该轴的平移的复合，
无不动点）。连同平移，即得平面等距变换的
完整分类。
\end{example}

\begin{remark}[平面等距变换一览]
归纳各种情形：恒等变换；平移（$\vec f =
\mathrm{id}$，除平凡情形外无不动点）；旋转
（线性部分为 $R_\theta$，$\theta \neq 0$：唯一中心）；
反射（线性部分为反射，不动点构成一条
直线）；滑动反射（线性部分同上，
无不动点）。除恒等变换外共四族，每族
仅由两项数据辨认：线性部分与不动点集。这正是
\cref{prop:b2:affine:fixedpoint}所述模式的
穷尽化。
\end{remark}

\begin{example}[当场捕获一个
滑动反射]\label{ex:b2:affine:glide}
设 $f(x, y) = (y + 1,\ x + 1)$。其线性部分 $(x, y)
\mapsto (y, x)$ 是关于对角线 $y = x$ 的反射，故
$1 \in \operatorname{Sp}\vec f$，而
\cref{prop:b2:affine:fixedpoint}无从给出结论。
不动点须同时满足 $x = y + 1$ 与 $y = x + 1$，
这不可能，故不存在不动点，$f$ 不是反射。
将其平方即可定出分类：
\[
f\bigl(f(x, y)\bigr) = f(y + 1,\ x + 1) = (x + 2,\ y + 2),
\]
即 $(2, 2)$ 的平移，故 $f$ 是\emph{滑动
反射}，其轴为直线 $y = x$（适当平移，
$M$ 与 $f(M)$ 的中点总位于
$y = x + {}$常数上（此处为 $y = x$），如取 $M = (0,
0) \mapsto (1,1)$ 可验证之），滑动向量 $(1, 1)$ 恰为 $f
\circ f$ 之半。与\cref{exo:b2:affine:6}相比，二者线性部分相同，
而常数项不同，所得即为一个真正的
反射：当\omterm{def:b2:reduction:eigen}{特征值} $1$ 出现时，常数项
决定一切。
\end{example}

\begin{example}[仿射递推即仿射
动力系统]\label{ex:b2:affine:recursion}
经典递推 $u_{n+1} = au_n + b$（$a \neq 1$）
迭代直线上的\omterm{def:b2:affine:subspace}{仿射映射} $f(x) = ax + b$，其
线性部分 $a$ 避开\omterm{def:b2:reduction:eigen}{特征值} $1$，故存在唯一
不动点 $\omega = \frac{b}{1-a}$，并在该点重新取原点
（即上述命题的一维情形）使
$f$ 化为乘以 $a$：
\[
u_{n+1} - \omega = a\,(u_n - \omega)
\qquad\Longrightarrow\qquad
u_n = \omega + a^n(u_0 - \omega) .
\]
对 $u_{n+1} = \frac{u_n}2 + 3$，有 $\omega = 6$，且 $u_n = 6 +
(u_0 - 6)2^{-n} \to 6$。此类递推在
\cref{ch:b2:series}中所教的方法——“减去不动点”——
正是\omterm{def:b2:affine:subspace}{仿射映射}在其不动点处的向量化；
当 $\abs a < 1$ 时的收敛，即
压缩现象，经\cref{ch:b2:metric}化为
巴拿赫不动点定理。一个想法，贯穿三章。
\end{example}

\begin{example}[求旋转的中心]\label{ex:b2:affine:rotationcenter}
设 $f(x, y) = (-y + 2,\ x)$。其线性部分为
$\varphi(x, y) = (-y, x)$，即转角为 $\frac\pi2$ 的旋转，
其谱 $\{\iu, -\iu\}$ 不含 $1$。由
\cref{prop:b2:affine:fixedpoint}，恰有一个
不动点。由 $x = -y + 2$ 与 $y = x$ 得 $x = 1$，$y = 1$，
故 $\Omega = (1, 1)$，且 $f$ \emph{即为}中心为
$(1, 1)$、转角 $\frac\pi2$ 的旋转。一般地，当 $1
\notin \operatorname{Sp}\vec f$ 时，对 $f$ 的分类仅需
一个线性方程组——几何完全体现于线性部分，
运算则完全在于确定中心。
\end{example}

\begin{remark}[仿射语言的后续应用]
\omterm{def:b2:affine:barycenter}{重心}与\omterm{def:b2:affine:subspace}{仿射映射}是后续几何各章的“语法”：
切线与切平面都是仿射对象
（\cref{ch:b2:curves,ch:b2:surfaces}），仿射变量替换使
面积与体积乘以
$\abs{\det \vec f}$（\cref{ch:b2:multint}），而期望
是以某一概率分布为权的\omterm{def:b2:affine:barycenter}{重心}，这正是
凸性支配琴生不等式的缘由
（\cref{ch:b2:randomvar}）。在本科三年级分册中，
同一套凸性词汇用于研究 $L^p$ \omterm{def:b2:nvs:norm}{范数}与
积分不等式。
\end{remark}

\begin{remark}[本册内的展望]
本章引出两条线索。其一是\emph{仿射}线索：
切线（\cref{ch:b2:curves}）与切平面
（\cref{ch:b2:surfaces}）是附着于
非线性对象上的仿射子空间，而曲面一章中
二次曲面的分类依赖于本章的中心方程
$A\Omega = -b$。另一条\emph{凸性}线索更长：凸性
半平面与圆盘的凸性支撑了
周末问题中的赫利理论；函数的凸性给出琴生
不等式（\cref{ch:b2:randomvar}）；而全书最后的定理——
分支过程的灭绝判别准则
（\cref{ch:b2:genfun}）——由一条凸曲线相对于
对角线的位置决定，这幅图既属于本章，
也同样属于概率论。\omterm{def:b2:affine:barycenter}{重心}
在此亦再次出现：期望即是
以概率为权的\omterm{def:b2:affine:barycenter}{重心}。
\end{remark}

\section{练习}

\begin{exercise}[$\star$]\label{exo:b2:affine:1}
在 $\R^3$ 中，下列集合是否为仿射子空间？请给出方向与
维数。$\{x + y + z = 1\}$；$\;\{x + y + z = 1,\ x - z =
3\}$；$\;\{x^2 + y^2 = 1\}$；给定相容方程组
$MX = B$ 的解集。
\end{exercise}

\begin{exercise}[$\star$]\label{exo:b2:affine:2}
证明\omterm{def:b2:affine:def}{仿射空间}中三个互异的点 $A, B, C$ 共线，
当且仅当 $C$ 是 $A$ 与 $B$ 的\omterm{def:b2:affine:barycenter}{重心}，又当且仅当向量
$\vect{AB}, \vect{AC}$ 线性相关。由此推出梅涅劳斯型的
权重计算：若 $C = \operatorname{bar}(A, 1 - t; B, t)$，
分别确定 $C$ 在 $t = \frac12$、$t = 2$、$t = -1$ 时的位置。
\end{exercise}

\begin{exercise}[$\star$]\label{exo:b2:affine:3}
设 $f$ 为 $\R^2$ 上由 $f(X) = MX + C$ 给出的\omterm{def:b2:affine:subspace}{仿射映射}，其中
$M = \frac12\begin{pmatrix} 1 & 1\\ 1 & 1\end{pmatrix}$，且 $C =
(1, 0)^{\mathsf T}$。求 $f$ 的像、
其不动点（若存在），以及 $f \circ f$。
\end{exercise}

\begin{exercise}[$\star\star$]\label{exo:b2:affine:4}
（结合律的应用）在三角形 $ABC$ 中，设 $I, J, K$ 分别将
$BC$、$CA$、$AB$ 按比例分割，使 $\vect{BI} = \frac13\vect{BC}$，
$\vect{CJ} = \frac13\vect{CA}$，$\vect{AK} = \frac13\vect{AB}$。
用\omterm{def:b2:affine:barycenter}{重心}表示 $I, J, K$，并求
$(I,1;J,1;K,1)$ 的\omterm{def:b2:affine:barycenter}{重心}。所得结果是什么？为何可预料？
\end{exercise}

\begin{exercise}[$\star\star$]\label{exo:b2:affine:5}
证明映射 $f \colon \R^n \to \R^n$ 保持\emph{中点}
（$f\bigl(\frac{A+B}{2}\bigr) = \frac{f(A) + f(B)}{2}$）且
\emph{\omterm{def:b2:metric:continuity}{连续}}，则它是\omterm{def:b2:affine:subspace}{仿射映射}。\emph{（证明向量映射 $u \mapsto
f(O + u) - f(O)$ 借助中点满足可加性，再
得 $\Q$-齐次，再由\omterm{def:b2:metric:continuity}{连续}性得 $\R$-齐次——
这与本科一年级分册中柯西函数方程所用的
稠密性策略相同；此处需重新推导所需步骤。）}
\end{exercise}

\begin{exercise}[$\star\star$]\label{exo:b2:affine:6}
对欧氏平面的\omterm{def:b2:affine:subspace}{仿射映射} $f(x, y) = (y + 1,\; x - 1)$ 进行分类：
求其线性部分、不动点与几何性质
（是反射还是滑反射？）。计算 $f \circ f$，并作出结论。
\end{exercise}

\begin{exercise}[$\star\star\star$]\label{exo:b2:affine:7}
（拉东）设 $A_1, \dots, A_{n+2}$ 为\omterm{def:b2:affine:def}{仿射空间}中的点，
该空间的维数为 $n$。证明可将它们分成两个不相交的
组，使两组的\omterm{def:b2:affine:convex}{凸包}相交。\emph{（仿照卡拉西奥多里的
证明，求权 $\mu_i$（不全为零），满足 $\sum\mu_i = 0$ 且
$\sum\mu_i\vect{OA_i} = 0$；将正权与负权分开，
并分别对两侧归一化。）}
\end{exercise}

\begin{exercise}[$\star\star\star$]\label{exo:b2:affine:8}
设 $f$ 为 $\R^n$ 上满足 $f \circ f = f$ 的仿射自同态。
证明 $f$ 是到仿射子空间
$\operatorname{Fix}(f) = \operatorname{im} f$ 上、沿方向
$\ker\vec f$ 的仿射投影，并证明反之，所有此类投影都是
幂等的。\emph{（先证明 $\operatorname{im} f$ 中
的点都是不动点。）}
\end{exercise}

\begin{exercise}[$\star$]\label{exo:b2:affine:9}
设 $G = \operatorname{bar}(A, 1;\ B, 2;\ C, 3)$ 为三角形 $ABC$ 中的点，
利用结合律，证明直线 $AG$ 与 $BC$
交于 $M = \operatorname{bar}(B, 2;\ C, 3)$，并确定 $G$ 位于
线段 $\intcc AM$ 上的位置；同样确定 $BG$
与 $CA$ 的交点。
\end{exercise}

\begin{exercise}[$\star\star$]\label{exo:b2:affine:10}
对 $\lambda \neq 0$，\emph{位似} $h_{\Omega,
\lambda}$ 是以 $\Omega$ 为不动点、线性部分为
$\lambda\,\mathrm{id}$ 的\omterm{def:b2:affine:subspace}{仿射映射}。证明复合
$h_{\Omega', \mu} \circ h_{\Omega, \lambda}$ 为
比值 $\lambda\mu$ 的位似（当 $\lambda\mu \neq 1$ 时），
而在 $\lambda\mu = 1$ 时为平移；当 $\lambda = \mu = -1$（两次中心对称）时，
求该平移向量。
\end{exercise}

\begin{exercise}[$\star\star$]\label{exo:b2:affine:11}
（梅涅劳斯）在三角形 $ABC$ 中，设 $A' \in (BC)$，$B' \in
(CA)$，$C' \in (AB)$，三点均异于顶点，
并由下式定义 $\alpha, \beta, \gamma$：$\vect{A'B} =
\alpha\,\vect{A'C}$，$\vect{B'C} = \beta\,\vect{B'A}$，
$\vect{C'A} = \gamma\,\vect{C'B}$。证明 $A', B', C'$ 共线
当且仅当 $\alpha\beta\gamma = 1$。\emph{（将每个点写成
两个顶点的\omterm{def:b2:affine:barycenter}{重心}；证明三点共线当且仅当它们
关于 $(A, B, C)$ 的\omterm{def:b2:affine:barycenter}{重心}坐标行构成
奇异的 $3 \times 3$ 矩阵。）}
\end{exercise}

\begin{exercise}[$\star\star\star$]\label{exo:b2:affine:12}
设 $K$ 为 $\R^n$ 的紧子集，证明其\omterm{def:b2:affine:convex}{凸包}是紧的。
\emph{（由\cref{thm:b2:affine:caratheodory}，
$\operatorname{conv}(K)$ 是紧集在\omterm{def:b2:metric:continuity}{连续}映射下的像。）}试举一个
$\R^2$ 中的例子，说明
\emph{闭}集的\omterm{def:b2:affine:convex}{凸包}未必是闭的。
\end{exercise}

\section{问题：从拉东到赫利、中心点与荣格
定理}

\begin{omfigure}
\begin{tikzpicture}[scale=0.9]
  \draw[omDef, thick] (0,0) -- (3,0) -- (0.9,2.6) -- cycle;
  \fill (0,0) circle (1.6pt) node[below left] {\small $A_1$};
  \fill (3,0) circle (1.6pt) node[below right] {\small $A_2$};
  \fill (0.9,2.6) circle (1.6pt) node[above] {\small $A_3$};
  \fill[omThm] (1.3,0.9) circle (1.8pt)
    node[below, omThm] {\small $A_4$};
\end{tikzpicture}
\qquad
\begin{tikzpicture}[scale=0.9]
  \draw[omDef, thick] (0,0) -- (3,0.3) -- (3.3,2.4) --
    (0.4,2.2) -- cycle;
  \draw[omThm, thick] (0,0) -- (3.3,2.4);
  \draw[omThm, thick] (3,0.3) -- (0.4,2.2);
  \fill (0,0) circle (1.6pt) node[below left] {\small $A_1$};
  \fill (3,0.3) circle (1.6pt) node[below right] {\small $A_2$};
  \fill (3.3,2.4) circle (1.6pt) node[above right] {\small $A_3$};
  \fill (0.4,2.2) circle (1.6pt) node[above left] {\small $A_4$};
  \fill[omProp] (1.75,1.27) circle (1.8pt);
\end{tikzpicture}

{\small 平面上一般位置四点的两种拉东型：
一点位于其余三点构成的三角形内部
（划分 $\{A_4\} \mid \{A_1, A_2, A_3\}$），
或处于凸位置，此时拉东点（橙色）
为两条对角线的交点。}
\end{omfigure}

\begin{problem}[{周末问题——赫利定理及其
两个应用}]\label{pb:b2:affine:1}
拉东引理（\cref{exo:b2:affine:7}）说明，任意 $n + 2$
个点（位于 $n$ 维\omterm{def:b2:affine:def}{仿射空间}中）总可分成两组，使两组的\omterm{def:b2:affine:convex}{凸包}相交。
本题将这一线性代数事实
串成一系列组合几何定理：
赫利交集定理、中心点定理（二维中位数）
以及荣格覆盖定理。全题中，平面取为
赋有通常欧几里得结构的 $\R^2$，
$\det$ 为在标准基下的\omterm{def:b2:linalg:det}{行列式}。
\index{赫利定理}\index{拉东
引理}\index{荣格定理}\index{中心点}

\medskip
\noindent\textbf{第一部分——\omterm{def:b2:affine:barycenter}{重心}坐标。}点
$A_0, \dots, A_k$ 称为\emph{仿射无关}，当
向量 $\vect{A_0A_1}, \dots, \vect{A_0A_k}$ 线性
无关。
\begin{enumerate}
  \item 证明仿射无关性不依赖于基点 $A_0$ 的选取，
        且它等价于：若两族权重各自之和为 $1$，
        并且它们定义出 $(A_0,
        \dots, A_k)$ 的同一\omterm{def:b2:affine:barycenter}{重心}，
        则两族权重相同。
  \item 设 $(A, B, C)$ 在平面中仿射无关，证明每一点 $M$
        都唯一对应一组三元数 $(\alpha, \beta, \gamma)$，
        满足 $\alpha + \beta +
        \gamma = 1$ 且 $M = \operatorname{bar}(A, \alpha; B,
        \beta; C, \gamma)$ ——称其为\emph{
        \omterm{def:b2:affine:barycenter}{重心}坐标}。
  \item 证明\omterm{def:b2:linalg:det}{行列式}公式
        \[
        \alpha = \frac{\det(\vect{MB},
        \vect{MC})}{\det(\vect{AB}, \vect{AC})},
        \qquad
        \beta = \frac{\det(\vect{MC},
        \vect{MA})}{\det(\vect{AB}, \vect{AC})},
        \qquad
        \gamma = \frac{\det(\vect{MA},
        \vect{MB})}{\det(\vect{AB}, \vect{AC})} :
        \]
        \omterm{def:b2:affine:barycenter}{重心}坐标即为有向面积之比。
  \item 直线 $BC$、$CA$、$AB$ 分别为坐标直线
        $\{\alpha = 0\}$, $\{\beta = 0\}$, $\{\gamma = 0\}$.
        证明 $M$ 位于闭三角形
        $\operatorname{conv}\{A, B, C\}$ 内，当且仅当 $\alpha, \beta,
        \gamma \geq 0$，并且三条直线将平面恰好分为七个区域，
        它们由
        $(\alpha, \beta, \gamma)$ 的符号加以区分（符号组 $(-, -,
        -)$ 不可能出现）。
  \item 设 $u \colon \R^2 \to \R$ 为\omterm{def:b2:affine:subspace}{仿射映射}（称为
        \emph{仿射型}）。证明 $u(M) = \alpha\,u(A) +
        \beta\,u(B) + \gamma\,u(C)$，非常值仿射型的水平集为直线，
        每条直线都以这种方式出现，
        并且闭半平面 $\{u
        \geq c\}$ 是凸的。
\end{enumerate}

\medskip
\noindent\textbf{第二部分——精细化的拉东划分。}若
$n + 2$ 个点构成 $\R^n$ 中的一个族，其中
任意 $n + 1$ 个点仿射无关，则称该族处于\emph{一般位置}。
\emph{仿射相关}是指 $(A_1, \dots, A_{n+2})$ 的
一族 $(\mu_i)$，满足 $\sum_i \mu_i = 0$ 且 $\sum_i
\mu_i\,\vect{OA_i} = 0$，对某一（从而对任一）原点 $O$ 成立。
\begin{enumerate}[resume]
  \item 求下列四点的一个非零仿射相关，
        $A_1 = (0,0)$，$A_2 = (3,0)$，$A_3 = (0,3)$，
        $A_4 = (1,1)$，并给出拉东划分与
        拉东点。
  \item 证明：对一般位置的点，仿射相关所构成的
        向量空间的维数恰为
        $1$，且任一非零仿射相关\emph{没有}
        为零的系数。
  \item 由此推出：一般位置的 $n + 2$ 个点的拉东划分
        是唯一的（至多交换两块的次序），
        每一块都是使
        $\mu_i$ 取某一固定符号的指标集。
  \item 对平面上一般位置的四点，证明下述二分性：
        划分的类型为
        $(1, 3)$，即一点位于其余三点所围三角形的内部；
        或为类型 $(2, 2)$：四点
        处于凸位置，且连接两对点的线段
        （即对角线）相交，交点即为
        拉东点。
  \item 按第 6 题，对单位正方形 $(0,0)$、
        $(1,0)$、$(1,1)$、$(0,1)$，求
        仿射相关、划分及拉东点。
\end{enumerate}

\medskip
\noindent\textbf{第三部分——平面上的赫利定理。}
\begin{enumerate}[resume]
  \item 设 $C_1, C_2, C_3, C_4$ 为 $\R^2$ 的凸子集，
        其中任意三个有公共点。取 $x_i \in
        \bigcap_{j \neq i} C_j$，并对
        $x_1, \dots, x_4$ 应用拉东引理，证明拉东点
        属于全部四个集合。\emph{（对每个 $k$，不含 $x_k$ 的
        那一块由 $C_k$ 中的点组成。）}
  \item （赫利）设 $C_1, \dots, C_m$（$m \geq 3$）为
        $\R^2$ 中的凸子集，其中任意三个相交。
        证明 $\bigcap_{i=1}^m C_i \neq \emptyset$，
        对 $m$ 用归纳法：将 $C_{m-1}$ 与 $C_m$ 换成
        $C_{m-1} \cap C_m$，并检验新族满足的假设，
        利用第 11 题。
  \item 三个反例，各对应一个假设：（a）三角形的
        三条闭边两两相交，
        但没有公共点（$3$ 不能降为
        $2$）；（b）四个集合 $S_i = \{x_1, \dots, x_4\}
        \setminus \{x_i\}$，对一般位置的四点，
        满足三点相交的假设，
        但不满足结论（凸性必不可少）；（c）
        闭半平面 $H_k = \intco{k}{+\infty} \times
        \R$，$k \in \N$，两两相交且三三相交，
        但 $\bigcap_k H_k = \emptyset$（无穷族
        需要紧性）。
  \item （紧赫利定理）设 $(K_i)_{i \in I}$ 为任一
        由 $\R^2$ 的紧凸子集组成的族，其中任意三个
        相交。利用第 12 题与
        博雷尔–勒贝格性质
        （见\cref{thm:b2:metric:borellebesgue}），证明
        $\bigcap_{i \in I} K_i \neq \emptyset$。
  \item （第一项收获）设 $S$ 为平面上的有限点集，
        $r > 0$。证明：若 $S$ 中任意三点
        都落在某个半径为 $r$ 的闭圆盘内，则 $S$
        落在某一个半径为 $r$ 的闭圆盘内。\emph{（对圆盘
        $\overline D(p, r)$，$p \in S$，应用赫利定理。）}
\end{enumerate}

\medskip
\noindent\textbf{第四部分——中心点定理。}\emph{中心点}是指
平面上有限集 $S$（共 $n$ 个点）的一点
$c$（不必属于 $S$），使得每个包含 $c$ 的闭半平面
都至少含有 $n/3$ 个
$S$ 中的点。
\begin{enumerate}[resume]
  \item （$1$ 维情形）对实数 $x_1 \leq \dots \leq x_n$，
        证明中位数 $c = x_{\lceil n/2 \rceil}$
        满足：包含 $c$ 的每个闭半直线
        都至少含有 $n/2$ 个 $x_i$。
  \item （计数引理）若 $A, B, C$ 为 $S$ 的子集，且
        $\abs A, \abs B, \abs C > \tfrac{2n}3$，证明 $A \cap
        B \cap C \neq \emptyset$。
  \item 令 $m = \floor{2n/3} + 1$，并设 $\mathcal F$ 为
        由\omterm{def:b2:affine:convex}{凸包}组成的有限族，
        $\operatorname{conv}(T)$，其中 $T \subseteq S$ 且 $\abs T =
        m$ 个点。证明 $\mathcal F$ 中任意三个成员都有
        公共点，并由赫利定理推出存在一点 $c$
        同时属于它们全部。
  \item 证明：此 $c$ 是 $S$ 的中心点，即
        \emph{中心点定理}。\emph{（若过
        $c$ 的闭半平面所含点少于 $n/3$ 个，则其开补集
        将含有一个集合 $T$，
        由 $m$ 个点组成，而 $\operatorname{conv}(T)$ 将
        不含 $c$。）}
  \item 最优性：设 $n = 3k$，在三个圆盘中各放 $k$ 个点，
        这些圆盘的半径为小量 $\varepsilon$，圆心
        位于一个大三角形的顶点处。证明：对
        \emph{每一个}平面上的点 $c$，都存在一个闭半平面
        包含 $c$，且至多含 $n/3$ 个
        $S$ 中的点，因此常数 $1/3$ 不能
        被改进。\emph{（从 $c$ 出发指向三个圆心的
        三个方向中必有两个所成角至多为
        $2\pi/3$。）}
\end{enumerate}

\medskip
\noindent\textbf{第五部分——荣格定理与综合。}
\begin{enumerate}[resume]
  \item （三角形引理）设 $P, Q, R$ 为三点，两两
        距离均 $\leq 1$。证明它们落在
        半径为 $1/\sqrt3$ 的闭圆盘内。\emph{（若某个
        角 $\geq \pi/2$，则取以最长边
        为直径的圆盘，利用中线公式
        $\norm{RM}^2 = \tfrac12\norm{RP}^2 +
        \tfrac12\norm{RQ}^2 - \tfrac14\norm{PQ}^2$；若
        三角形为锐角时，则估计外接圆半径 $a/(2\sin
        \widehat A)$，并用其最大角，该角位于
        $\intco{\pi/3}{\pi/2}$ 内。）}
  \item （荣格）由此推出：平面上直径
        $\leq 1$ 的每个紧子集都含于一个
        半径为 $1/\sqrt3$ 的闭圆盘中。
  \item 最优性：对正三角形 $A_1A_2A_3$，
        边长为 $1$，\omterm{def:b2:affine:barycenter}{重心}为 $G$，证明莱布尼茨
        恒等式 $\sum_i \norm{\vect{OA_i}}^2 =
        3\norm{\vect{OG}}^2 + \sum_i \norm{\vect{GA_i}}^2$
        对每一点 $O$ 成立，并由此推出：包含三顶点的任一圆盘
        的半径满足 $\geq
        1/\sqrt3$，且仅对外接圆盘取等号。
  \item （赫利，$\R^n$ 情形）陈述并证明
        赫利定理：若 $\R^n$ 中的有限个\omterm{def:b2:affine:convex}{凸集}
        满足任意 $n + 1$ 个都相交，则它们全体相交。
        \emph{（拉东引理
        见\cref{exo:b2:affine:7}，可处理 $n + 2$ 个集合；再
        仿照第 12 题作归纳。）}
  \item 综合。写出下面的链条
        \[
        \text{仿射相关} \Rightarrow \text{拉东}
        \Rightarrow \text{赫利} \Rightarrow
        \text{中心点与荣格},
        \]
        每项各用一句话说明：线性代数在何处进入，
        权重的符号在何处进入，
        凸性在何处进入，以及哪一步用到了
        平面的维数。常数 $3$（赫利定理中的）、
        $1/3$（中心点）与 $1/\sqrt3$（荣格）
        在 $\R^n$ 中分别变为何？（仅陈述，不证明。）
\end{enumerate}
\end{problem}
